
\section{Data Sets}\fbeg{}\begin{frame}{\ft{Data Sets}}

{\fontsize{20}{22}\selectfont

\hspace*{-2pt}\begin{ftxt}[0.95\textwidth]

Accessibility has multiple dimensions.  Starting from a published 
manuscript, there are usually several steps required to download, 
deserialize, and study data sets.  With FAIR data, these steps 
can hopefully be accomplished without relying on 
a lot of external software or manual processing.  
Ideally, computer code needed to unpack and visualize 
data contents are included as part of a data set.\\\vspace{-13pt}

FAIR data sets should not rely on users having access to 
external software, especially software that is not 
free and open-source.  Researchers should be able 
to assess scientific work without having the same 
applications or tools that the original scientists employed.
For most data sets, raw data needs special deserialization 
and visualization tools.  For FAIR/Accessibility, 
XML files can't just be viewed as DOM trees in web 
browsers, and CSV can't just be opened as spreadsheets.


\end{ftxt}

}

\end{frame}

